\section{Discussion}
The central lesson of the current artifact is that a prompt-selection project needs a stopping rule as much as it needs a selection model. Without Gate F, the natural engineering response to a small gap would be to add more actions, more features, or more training variants until one branch appears favorable. AdaPromptSet instead treats the oracle gap as an upstream feasibility check. The per-query oracle is allowed to be strong only as an analysis object; it must expose enough headroom before any learned selector can be justified. In the present MoNuSAC pilot, that headroom is too small.

This does not make the experiment empty. The fixed action \OracleAnalysisBestFixed{} is a useful reference point because it combines top candidate positives with negative evidence under a simple call budget. The fact that it is close to the bank oracle suggests that much of the available value in the current candidate representation is already captured by this fixed structured rule. At the same time, the action summaries show that oracle wins are not perfectly concentrated in a single action. The correct reading is therefore balanced: the bank contains some query-specific variation, but not enough measured benefit to justify training a policy on MoNuSAC.

The representation diagnostics sharpen this interpretation. Low component coverage indicates that the candidate generator often fails to touch enough target components. If the evidence does not cover the objects, then downstream action ranking has limited room to succeed. The top-two margin diagnostic shows that actions are separated in some episodes, which prevents the paper from claiming that all actions are equivalent. The SAM sensitivity diagnostic shows no clear advantage for higher-call grouped actions in this pilot. Together, these diagnostics suggest that the next experiment should be designed around candidate coverage and domain-specific evidence quality, not around a larger policy head.

A second lesson is that external data must be treated as a protocol decision rather than a convenient rescue. GlaS is available locally and has a clean morphology-shift role, but its present validation evidence has only one strict episode. DDR has enough paired labels for infrastructure checks, but the stage-one union task is not the formal four-intent lesion protocol. IDRID is recorded as read-only because prior observations could bias future choices. CPM-17 remains deferred because the local assets do not meet the expected complete public-version manifest. These details make the paper less dramatic, but more useful: they tell the next experiment exactly where it can and cannot draw evidence.

\subsection{What Evidence Would Change the Conclusion}
The current conclusion would change only through evidence that is both predeclared and independent of the failed MoNuSAC rescue path. Three conditions are especially important. First, a future action bank must improve candidate coverage before action selection is evaluated. The present post-failure audit points to low component coverage, so another selector trained on the same weak candidate set would mostly learn from a narrow set of already observed successes and failures. A stronger continuation would therefore register the candidate generator, coverage metric, and acceptance threshold before computing any new oracle labels.

Second, action heterogeneity must be large enough to matter after matched-budget accounting. The current oracle wins are split across actions, but the mean fixed-oracle Dice gap is small. A future gate should require not only a larger mean gap, but also enough query units to estimate whether the gap is stable across cases rather than driven by one or two episodes. This protects the project from promoting an action selector when the observed variation is real but too small, or when the apparent gain is a sampling artifact. The paired sign-flip result in the present artifact is a useful warning: even a favorable mean can be statistically unconvincing when the number of independent units is tiny.

Third, external-domain evidence must be collected under a fresh role assignment. The next domain should not be selected because it makes the current paper easier to write; it should be selected because its manifest, reference pools, masks, and split logic are complete enough to support the gate before any method choice is tuned. GlaS is currently the most natural morphology-shift continuation because its local assets are complete, but the one-episode smoke is only a wiring check. DDR can stress image scale and lesion morphology, but the available stage-one union protocol is a different task. These distinctions keep follow-up experiments from borrowing credibility from this MoNuSAC pilot.

\subsection{Gate Extensions After a Failed Node}
A practical gate sequence should also say what happens when a node fails for a specific reason. If coverage fails, the allowed work is a candidate-generator diagnostic, such as component coverage by structure size, point precision against a random baseline, or a reference-quality sensitivity table. If action heterogeneity fails, the allowed work is an action-bank analysis that reports fixed-action rankings, oracle win counts, near ties, and matched-budget deltas without training a selector. If statistical support fails, the allowed work is a larger predeclared validation run or a patient-cluster audit, not a stronger wording of the same small pilot. If an external manifest fails, the allowed work is dataset repair and split registration, not an improvised positive result.

This failure-routing view is the main gate-mechanism refinement suggested by the current paper. The gates should not be only pass/fail labels. They should map each failure mode to a narrow set of legal next actions and a narrow set of prohibited shortcuts. In the current state, the legal next actions are fixed-rule analysis, representation diagnostics, limited candidate diagnostics, GlaS readiness expansion, DDR infrastructure smoke, and CPM-17 asset repair. The prohibited shortcuts remain MoNuSAC policy training, post-hoc MoNuSAC action-bank redesign, and query-time action-bank search. This makes the project easier to continue because a failed node does not leave the researcher guessing; it opens only those experiments that can plausibly repair the missing evidence.

Table~\ref{tab:gate-routing} turns this rule into an executable review checklist. The table is intentionally conservative: each failure mode has a diagnostic that can be run without changing the claim, a minimum record needed before promotion, and a shortcut that would invalidate the current interpretation. This is the main mechanism by which AdaPromptSet can use additional local datasets without letting them become ad hoc evidence.
\begin{table*}[t]
\centering
\small
\caption{Failure-to-experiment routing used after the MoNuSAC Gate-F pivot. Allowed continuations preserve the oracle-analysis-only claim unless a future predeclared gate passes with enough independent units.}
\label{tab:gate-routing}
\begin{tabular}{@{}p{0.18\textwidth}p{0.30\textwidth}p{0.26\textwidth}p{0.18\textwidth}@{}}
\toprule
Failed evidence node & Allowed continuation & Minimum promotion record & Prohibited shortcut \\
\midrule
Candidate coverage & Component coverage by target size, point precision against random baseline, and reference-quality sensitivity. & Predeclared coverage threshold plus frozen train/validation split before utility labels are regenerated. & Retuning prompts on MoNuSAC validation cases after seeing oracle wins. \\
Action heterogeneity & Fixed-action ranking, oracle win counts, near-tie fraction, and matched-budget deltas. & Larger oracle gap with enough query or patient units to make the gap interpretable. & Training a selector from a small or near-tie utility matrix. \\
Statistical support & Query-level paired analysis, patient-cluster audit, and bootstrap only when clusters exist. & Independent case or patient count sufficient for the declared interval or sign test. & Rewording a small pilot as a formal positive result. \\
External manifest & Dataset repair, split registration, read-only marking, or infrastructure smoke according to dataset role. & Complete image-mask pairs, non-leaky references, and a fresh domain gate. & Choosing a dataset because it improves the story after MoNuSAC fails. \\
Cost accounting & Same-hardware runtime, SAM-call counts, prompt counts, and GPU-hour ledger rows. & Quality and wall-clock evidence reported together under one frozen protocol. & Reporting speed or cost separately from the quality gate. \\
\bottomrule
\end{tabular}
\end{table*}

\section{Reproducibility and Audit Trail}
The current repository is organized as a self-contained MVP package with source code, configs, manifests, reports, and paper assets. The handoff report records the phase status, gate outputs, fixed-oracle summary, representation summary, new-domain readiness, bootstrap status, compute ledger, run metadata, PDF hash, and reproduction commands. The paper is generated from scripts rather than edited as a standalone artifact; submission mode rejects simulated placeholders, and the compiled PDF metadata is audited after build.

The reporting chain is intentionally redundant. \texttt{reports/final\_project\_report.json} is the machine-readable summary. \texttt{reports/scientific\_claim\_audit.json} checks that reader-facing files do not contain prohibited claims and that required boundary evidence exists. \texttt{reports/paper\_section\_audit.json} checks that the abstract, related work, method, experiments, results, limitations, and conclusion remain aligned with the oracle-analysis-only state. \texttt{reports/pdf\_metadata\_audit.json} checks the visible title and PDF info fields. This makes the paper rebuildable and keeps future edits from quietly changing the claim boundary.

The current PDF is not a submission-ready positive-result paper. It is an internally reviewable oracle-analysis paper that documents why the adaptive-policy branch stopped. That distinction is visible in the title metadata, the abstract, the method figure, the action-bank figure, the oracle captions, and the final report. A reviewer can reproduce the path from data audit to gate decision without treating the oracle as an automatic method. This audit trail is the main deliverable at the present state.

\section{Limitations}
The current action bank is an MVP subset and therefore an action bank upper bound rather than a complete prompt design search. Reference quality sensitivity remains unquantified in the current artifacts. Frozen matcher errors cannot be recovered by the current selector because no online SAM-response refinement or continuous coordinate generation is allowed. Utility generation cost remains part of the evidence: running the action bank consumes SAM calls, runtime, and GPU-hours even when the final deployable rule is fixed. GlaS has only one strict validation episode in the present smoke, and no formal frozen test has been run. No policy, transfer, efficiency, or leaderboard-level performance claim is supported by the current artifacts.

\subsection{What the Pivot Does Not Prove}
The failed gate does not prove that adaptive prompting is impossible. It proves only that the current bank, candidate evidence, and MoNuSAC MVP split do not justify the next policy-training phase. The low component coverage suggests that some target structures are not well represented by the present candidates. Conversely, the lack of higher-call gain suggests that simply adding grouped prompts is not a guaranteed fix. Both observations narrow the next experimental question, but neither is a deployable conclusion.

The paper also does not claim that \OracleAnalysisBestFixed{} is a final deterministic method. Its validation performance is near the action-bank oracle in the pilot, but deterministic promotion would require stronger stability evidence, broader query units, and a fair comparison to fully implemented baselines such as SynPo-style and GF-SAM-style variants with all costs included. Since those rows are not complete formal comparators in the current artifact set, they are not used to create a stronger title or abstract claim.

\subsection{Reporting Boundaries}
All oracle rows are analysis upper bounds. Privileged prompts, when present in diagnostics, are API and coordinate checks rather than automatic candidates. Multi-reference episodes are repeated measurements. The official GlaS test partitions are preserved; the one-episode smoke cannot be promoted to a gate pass. These boundaries are deliberately repetitive across the paper and reports because they are the difference between a useful negative result and an overfit narrative.
